\documentclass[12pt,titlepage,french]{article}
\usepackage{babel}
\usepackage{graphicx}
\usepackage[margin=2.5cm]{geometry}

\usepackage[hidelinks]{hyperref}
\usepackage{tabularx}
\usepackage{float}
\usepackage[utf8]{inputenc}
\usepackage[T1]{fontenc}
\pagestyle{plain}

\usepackage{booktabs,makecell,tabu}
\usepackage{comment}
\renewcommand\theadfont{\bfseries}

\linespread{1.5}

\newcounter{firstbib}

\begin{document}
%\renewcommand{\thesection}{\arabic{section}} % utilisé pour spécifier la numérotation des sections

\begin{titlepage}
\newcommand{\HRule}{\rule{\linewidth}{0.5mm}}
\center

  \includegraphics[width=0.45\textwidth]{../../ressources/img_logos/logo_polytech.png}\\[1cm]

  \includegraphics[width=0.45\textwidth]{../../ressources/img_logos/logo_taglabs.png}


\HRule \\[0.4cm]
{ \huge \bfseries Rapport itération 5\\[0.15cm] }
Classification colorimétrique de nuages de points 3D\\
Version 1.0\\
Le \today \\
\HRule \\[1.5cm]
Ronan Collier,
Mathieu Letrone,
Tri-Thien Truong
\\[1cm]
\end{titlepage}

\tableofcontents % table des matières
\newpage
\listoffigures  % table des figures
\newpage

\section{Rappel des objectifs de l'itération}
Suite à l'itération où nous avions amélioré notre solution, nous avons voulu élargir nos types de filtrages, en utilisant d'autres espaces colorimétriques. De plus, nous avions encore des tâches en cours, que nous devions avancer/finaliser.

Durant l'itération précédente, nous avons pu implémenter un nouveau filtre, en utilisant l'espace colorimétrique HSV. Il nous fallait maintenant avoir des retours clients pour que l'on puisse corriger les différents bugs. De plus, des tâches majeurs étaient encore en cours.

Les tâches que nous nous sommes fixées sont les suivantes :

\begin{itemize}
    \item Restreindre l'utilisateur lors de la saisie du filtrage RGB (borne inférieure et supérieure)
    \item Amélioration des plages du filtrage HSV + test de performance
    \item Résoudre les divers problèmes techniques qui empêchent le fonctionnement de l'algorithme de segmentation
    \item Résoudre problèmes et débuggage du toon mapping
    \item Comparer les performances avec un autre algorithme de toon mapping
\end{itemize}

Nous avions aussi une tâche qui s'est ajoutée pendant l'itération, qui était d'exporter le plugin pour pouvoir l'utiliser sur toutes les machines, en utilisant CloudCompare. Nous avions aussi eu des retours utilisateurs, afin d'améliorer notre plugin.

\section{Production / réalisation durant l'itération}

Nous développerons ici chaque objectif que nous nous sommes fixé pour cette itération.

\subsection{Restreindre l'utilisateur lors de la saisie du filtrage RGB}

Lorsque nous avons décidé cette tâche, nous voulions améliorer le filtrage RGB côté utilisateur, afin de mieux choisir les bornes minimum et maximum. \newline

En effet, l'utilisation de l'espace colorimétrique RGB limite les bornes que l'on peut utiliser, afin de sélectionner les points que l'on veut. Pour que ce filtrage fonctionne correctement, nous sommes obligé de sélectionner deux points qui ont la même gamme de couleur. Il faut ensuite choisir le point le plus "sombre", c'est-à-dire avec les valeurs RGB les plus faibles, puis le point le plus clair. \newline

Plusieurs problèmes peuvent se poser. Par exemple, si l'utilisateur décide de prendre des points dans des gammes de couleurs différentes, ce qui donne un résultat non satisfaisant. Il peut aussi avoir un problème lorsque l'utilisateur inverse la sélection des points, ce qui inverse les bornes. Il n'y aurait donc aucun point de sélectionné. \newline

Une solution peut consister à inverser les valeurs, si l'utilisateur se trompe. Pour cela, il faut être sûr que l'utilisateur a bien choisi deux points de la même gamme de couleur, puis vérifier que les trois composantes du deuxième point sont inférieurs au premier point, afin de les inverser. \newline

Une question reste en suspend, qui est de savoir si nous réalisons le traitement directement sur l'interface (l'utilisateur verra donc directement le changement des bornes), ou si le traitement se fera au niveau en pré-traitement du parcours de points, et l'utilisateur aura juste à choisir ses deux points.

\subsection{Amélioration des plages du filtrage HSV}

\subsection{Algorithme de segmentation}

\subsection{Résoudre problèmes et débuggage du toon mapping}

\subsection{Comparer les performances avec un autre algorithme de toon mapping}

\subsection{Exportation du plugin}

\section{Risques éliminés durant l'itération}


\section{Feedback}

\section{Commentaires sur l'itération}

Cette section va présenter nos ressentis sur notre itération. Cela peut correspondre à la façon dont nous avons pu gérer la charge de travail que nous avions prévu en début d'itération, des potentiels imprévus, points positifs/négatifs, et autres.

\subsection{Commentaires sur l'itération de façon générale}

\subsection{Commentaires sur les méthodes de travail/changements de méthode}


\section{Objectifs de la prochaine itération}


\section{Résumé}
\subsection{Tâches principales réalisées dans l'itération}

\subsection{Tâches principales à réaliser pour la prochaine itération}

\begin{thebibliography}{3}

\end{thebibliography}
\end{document}
